\section{Introduction}

Root canal treatment is an essential procedure for preserving natural teeth, and accurate determination of the working length is key to its success: extrusion of filling material through the apical foramen markedly increases the risk of subsequent tooth loss~\cite{Ng2011}, whereas an underfilled canal is a factor associated with poor retreatment outcomes~\cite{Sabeti2024}. In clinical practice, the working length is usually determined during treatment with an electronic apex locator supplemented by intraoperative radiographs~\cite{Kaur2024}.

Preoperative periapical radiographs are routinely acquired before root canal treatment. Estimating the root length from these images before treatment could help clinicians anticipate the working-length range, preselect instrument lengths, and identify anatomical variations early. However, length measurement on periapical radiographs faces two fundamental difficulties. First, clinically exported images often carry no pixel-to-millimetre scale, and different export settings lead to widely varying pixel scales between images. Second, a periapical radiograph is a two-dimensional projection of a three-dimensional tooth, so projection angle and tooth inclination elongate or foreshorten the image, making the ratio between image length and true length vary from image to image.

Advances in computer vision and medical image analysis have enabled deep learning to be widely applied to detection, segmentation, and geometric measurement in dental radiographs~\cite{Aminoshariae2021,Fontenele2025}. However, most studies involving length measurement assume a known pixel scale or rely on a physical calibration object in the image; these assumptions rarely hold for retrospective data in which only de-identified exported images are available. Moreover, exposure conditions and the overlap of soft tissue and alveolar bone vary considerably between periapical radiographs, and the root apex often has low contrast against the surrounding bone. Segmenting the unprocessed image directly may therefore localise the apical boundary unstably and amplify the length-estimation error.

To address these challenges, this study proposes a calibration-free three-stage pipeline. First, a lightweight YOLO11n detector localises and extracts the target tooth from the full radiograph, restricting subsequent analysis to the clinically relevant region. Second, a morphological top-hat transformation enhances local contrast at the root, root canal, and alveolar-bone boundaries while suppressing large-scale exposure non-uniformity, and YOLO11n-seg produces the tooth mask. Third, the pixel length is measured along the mask's principal axis, normalised by the radiograph's long side to remove export-scale differences, and converted to millimetres by Gaussian process regression (GPR) with patient age, followed by a fixed safety offset. Each image is processed independently with fixed preprocessing and measurement parameters; the method should therefore be regarded as an integrated detection--segmentation--regression pipeline rather than a single end-to-end trainable network.

The main contributions of this work are summarised as follows:
\begin{itemize}
  \item A YOLO11n detector extracts the target tooth from full periapical radiographs, and aspect-ratio-preserving scaling with zero padding standardises the input size without distorting root geometry.
  \item A morphological top-hat transformation is used for contrast enhancement before segmentation; under identical segmentation and regression settings, it yields a lower length-estimation error than CLAHE and unprocessed images.
  \item A long-side normalisation that uses only the image size removes export-scale differences without a pixel-to-millimetre scale or physical calibration object.
  \item GPR with patient age converts the normalised length to millimetres, and a fixed safety offset biases the estimates towards the short side for clinical safety.
  \item Using 67 anterior teeth from different patients as an independent test set, preprocessing methods and regression models are compared one factor at a time to assess accuracy and clinical safety.
\end{itemize}

The remainder of this paper is organised as follows. Section~2 reviews related work on deep learning for dental radiographs and working-length estimation. Section~3 describes the proposed pipeline, including dataset preparation, tooth instance segmentation, and filling-length estimation. Section~4 presents the implementation environment, training settings, and evaluation metrics. Section~5 reports the segmentation performance, the comparisons of preprocessing and regression models, and the input-feature analysis. Section~6 discusses the main findings and limitations, and Section~7 concludes the paper and outlines future work.
